\documentclass[10pt]{article}
\usepackage[margin=.85in]{geometry}
\usepackage[T1]{fontenc}
\usepackage{lmodern,amsmath,amssymb,amsthm,mathtools,booktabs,microtype}
\usepackage[colorlinks=true,linkcolor=blue]{hyperref}
\newtheorem{theorem}{Theorem}[section]
\newtheorem{lemma}[theorem]{Lemma}
\newtheorem{corollary}[theorem]{Corollary}
\newtheorem{proposition}[theorem]{Proposition}
\newcommand{\E}{\mathbb E}
\newcommand{\R}{\mathbb R}
\newcommand{\ip}[2]{\langle #1,#2\rangle_n}
\title{Signed residual dynamics and a local noisy crossing certificate}
\author{Research continuation; initialization-to-entry remains open}
\date{}
\begin{document}
\maketitle
\begin{abstract}
The repository's strong-block statistic is the derivative of a smooth
parameter functional. Its signed residual representation separates strong-
and weak-cluster forcing without discarding complement neurons or training
gate changes. We prove local comparison and enclosure lemmas and give a
computer-assisted, fixed-dataset reversal certificate from a small
neighborhood of a learned canonical state. The full OOD error decreases
and rebounds on an explicit sector; weak-cluster forcing dominates at the
late endpoint. The certificate is conditional on reaching its entry
neighborhood. The canonical initialized reachability inequality and the
high-probability iid theorem remain open.
\end{abstract}

\section{Conventions and the signed weak source}
Let $f_\theta(x)=\sum_{i=1}^h w_i(u_i^Tx)_+$ in $\R^2$,
$e(x)=x-f_\theta(x)$, and $\mathcal L=\frac12\E_n\|e\|^2$.
Use the Euclidean parameter norm and the empirical output inner product
$\ip{v}{z}=\E_n v(x)^Tz(x)$. Throughout, the flow is compatible gradient
flow with the ReLU chain rule and loss dissipation; all differential identities
are almost everywhere. Training selectors are the actual selectors.
The cohort $I$ is fixed, including when it was chosen by an offline ID rule.
Set
\[
 M=\sum_{i\in I}w_i u_i^T,\qquad v=(r,-1)^T,\qquad
 Y=e_1^TMv-r,\qquad \ell=(Y,1)^T,\quad 0<r<1.
\]
The repository statistic, without inserting geometric $k$ factors, is
\begin{equation}\label{eq:q}
 Q_G=Y(r\dot m_{11}-\dot m_{12})+r\dot m_{21}-\dot m_{22}
     =\ell^T\dot Mv.
\end{equation}
Here $v$ is an auxiliary vector, not the unit OOD probe
$\xi=(r,-\sqrt{1-r^2})$.

Write $u_i=(a_i,b_i)$ and $w_i=(c_i,d_i)$. On the weak cluster, define
\begin{align*}
 F_I(x)&=\sum_{i\in I}c_i(u_i^Tx)_+,& F_O(x)&=f_1(x)-F_I(x),\\
 H_I(x)&=x_2\sum_{i\in I}\alpha_i(x)c_id_i,
 &L_I(x)&=x_2\sum_{i\in I}\alpha_i(x)c_i(d_i-b_i),\\
 Z_I(x)&=-x_1\sum_{i\in I}\alpha_i(x)c_ia_i.
\end{align*}
The weak source singled out in the supplied proof is
$\mathcal S_2=\frac12\E_{2,n}[(f_1-x_1)H_I]$.

\begin{lemma}[Exact square and signed corrections]\label{lem:source}
One has $H_I=F_I+L_I+Z_I$ and
\begin{equation}\label{eq:square}
\begin{split}
\mathcal S_2={}&\tfrac12\|F_I\|_{2,n}^2
 +\tfrac12\E_{2,n}[F_OF_I]-\tfrac12\E_{2,n}[x_1F_I]\\
 &+\tfrac12\E_{2,n}[(f_1-x_1)L_I]
 +\tfrac12\E_{2,n}[(f_1-x_1)Z_I].
\end{split}
\end{equation}
No sign is asserted for the four corrections.
\end{lemma}
\begin{proof}
Selector compatibility gives $F_I=\sum_I\alpha_i c_i(a_ix_1+b_ix_2)$.
Subtract this from $H_I$ and expand $f_1=F_I+F_O$ in its inner product
with $F_I$. These operations are pointwise finite-sum identities and do not
freeze a training gate.
\end{proof}
Thus approximate input--output alignment is not innocuous unless its
\emph{signed contribution at the scale of the source} is controlled. Small
noise in coordinates also need not give a smaller-order target term.

\section{A smooth potential and exact residual dynamics}
Let $J_t$ be the selected Jacobian of the training output map from parameters
to the empirical output space. The flow and residual equation are
\begin{equation}\label{eq:residual}
 \dot\theta=J_t^*e_t,\qquad \dot e_t=-K_te_t,\qquad K_t=J_tJ_t^*\succeq0.
\end{equation}
The adjoint includes empirical averaging. Matrix calculations below can instead
use outputs and residuals divided by $\sqrt N$ and ordinary transposes.

\begin{proposition}[The empirical statistic as a residual functional]\label{prop:potential}
Define the smooth polynomial in the parameters
\[
 \Phi_I(\theta)=\tfrac12(e_1^TMv-r)^2+e_2^TMv,\qquad
 g_t=\nabla_\theta\Phi_I(\theta_t),\qquad C_t=J_tg_t.
\]
Then $Q_G=\dot\Phi_I=\ip{C_t}{e_t}$. Its gradient is explicitly
\[
 (g_t)_{w_i}=\mathbf1_{i\in I}\ell(u_i^Tv),\qquad
 (g_t)_{u_i}=\mathbf1_{i\in I}v(w_i^T\ell).
\]
Equivalently,
\[
 C_t(x)=\sum_{i\in I}\alpha_i(x)
 \big[(u_i^Tx)(u_i^Tv)\ell+(x^Tv)(w_i^T\ell)w_i\big].
\]
All neurons, including those outside $I$, enter $e_t,J_t,K_t$.
\end{proposition}
\begin{proof}
Differentiate $\Phi_I$ with respect to $M$ to obtain $\ell v^T$, and then
with respect to the two neuron vectors to obtain $g_t$. The product rule gives
$\dot\Phi_I=\ell^T\dot Mv=Q_G$. Substitute the flow and take the adjoint.
Applying the network Jacobian to $g_t$ gives the last formula. Equation
\eqref{eq:residual} follows from the absolutely continuous ReLU chain rule,
including at preactivations that vanish on a set of positive time measure.
\end{proof}
The positivity of $K_t$ does not determine the sign of $\ip{C_t}{e_t}$.
It does, however, prevent exponential amplification in a comparison of
\emph{training residuals}. This is different from a parameter-space
Gr\"onwall comparison or a global frozen-kernel training assumption.

\begin{theorem}[Contractive finite-window comparison]\label{thm:comparison}
Fix an entry time $a$ and a horizon $H>0$. Write
$e_a=e$, $J_a=J$, $K_a=K$, $g_a=g$, $C_a=C$, and define
\[
 \bar e(\tau)=e^{-K\tau}e,\qquad
 \Psi(\tau)=\ip{C}{e^{-K\tau}e},\qquad 0\le\tau\le H.
\]
For almost every $\tau$, the exact discrepancy is
\begin{equation}\label{eq:directional}
\begin{split}
 Q_G(a+\tau)-\Psi(\tau)={}&\ip{C_{a+\tau}-C}{e_{a+\tau}}\\
 &-\int_0^\tau
 \ip{e^{-K(\tau-s)}C}{(K_{a+s}-K)e_{a+s}}\,ds.
\end{split}
\end{equation}
If $\|C_{a+s}-C\|_n\le d_C(s)$ and
$\|K_{a+s}-K\|_{\rm op}\le d_K(s)$ almost everywhere, then
\begin{equation}\label{eq:rho}
 |Q_G(a+\tau)-\Psi(\tau)|\le
 \|e\|_n\left[d_C(\tau)+\int_0^\tau
 d_K(s)\|e^{-K(\tau-s)}C\|_n\,ds\right].
\end{equation}
In particular uniform bounds $d_C,d_K$ give the uniform error
$\rho=\|e\|_n(d_C+H d_K\|C\|_n)$. There is no exponential growth factor.
\end{theorem}
\begin{proof}
Dissipation or \eqref{eq:residual} gives
$\|e_t\|_n\le\|e\|_n$ for $t\ge a$. Variation of constants relative to
$K$ gives
\[
 e_{a+\tau}-e^{-K\tau}e
 =-\int_0^\tau e^{-K(\tau-s)}(K_{a+s}-K)e_{a+s}\,ds.
\]
Subtract the two residual functionals and use self-adjointness of $K$ to
move its exponential onto $C$. This proves \eqref{eq:directional} with its
signs intact. Cauchy--Schwarz proves \eqref{eq:rho}. Since $K\succeq0$,
$\|e^{-Ku}\|\le1$ for $u\ge0$, proving the uniform bound. The integrals
only require measurable bounded Jacobians on the compact parameter path;
training-gate jumps are allowed.
\end{proof}
For a late positive witness, a lower bound on the signed expression
\eqref{eq:directional} suffices. Requiring its absolute value to be small
is stronger and can lose helpful drift.


\begin{corollary}[Directional budget with signed endpoint drift]\label{cor:directional}
For a fixed $\tau$, put $b_s=J_{a+s}^*e_{a+s}$,
$\Delta J_s=J_{a+s}-J_a$, and $z_s=e^{-K_a(\tau-s)}C_a$. Define
\begin{align*}
 P_\tau&=(g_{a+\tau}-g_a)^Tb_\tau
                  +g_a^T\Delta J_\tau^*e_{a+\tau},\\
 B_\tau&=\int_0^\tau\left[
       \|\Delta J_s^*z_s\|\,\|b_s\|
       +\|J_a^*z_s\|\,\|\Delta J_s^*e_{a+s}\|\right]ds.
\end{align*}
Then almost everywhere in $\tau$,
\begin{equation}\label{eq:directional-budget}
 Q_G(a+\tau)\ge\Psi(\tau)+P_\tau-B_\tau.
\end{equation}
In particular $P_\tau\ge-p$ and $B_\tau\le B$ imply a positive witness
whenever $\Psi(\tau)>p+B$. No upper bound on helpful positive endpoint
drift is required. These bounds on an interval give an increasing interval
for $\Phi_I$; a single time is not by itself such an interval.
\end{corollary}
\begin{proof}
The first term of \eqref{eq:directional} equals $P_\tau$, since
$J_tg_t-J_ag_a=J_t(g_t-g_a)+(J_t-J_a)g_a$. Moreover
\[
 K_{a+s}-K_a=\Delta J_sJ_{a+s}^*+J_a\Delta J_s^*.
\]
Thus its integrand is exactly
$(\Delta J_s^*z_s)^Tb_s+(J_a^*z_s)^T\Delta J_s^*e_{a+s}$.
Cauchy--Schwarz bounds this above by the integrand of $B_\tau$.
This proves the lower bound, retaining the sign of $P_\tau$.
\end{proof}
The first integral can alternatively be bounded by
\[
 \sqrt{\mathcal L(a)-\mathcal L(a+\tau)}
       \left(\int_0^\tau\|\Delta J_s^*z_s\|^2ds\right)^{1/2},
\]
using loss dissipation. This need not be sharper than keeping the pointwise
product. Neither form requires an operator-norm bound over all residual
directions.

\section{Spectral crossing and full-error transfer}
The comparison function is calculated in parameter dimension rather than
forming a $2N$-dimensional kernel. Put $B=J^*J$, $b=J^*e$. If
$Bz_j=\lambda_jz_j$ is an orthonormal eigendecomposition, then
\begin{equation}\label{eq:spectrum}
 \Psi(\tau)=g^Te^{-B\tau}b
     =\sum_j c_je^{-\lambda_j\tau},\qquad
 c_j=(g^Tz_j)(z_j^Tb),\quad \lambda_j\ge0.
\end{equation}
This follows from $J^*e^{-JJ^*\tau}=e^{-J^*J\tau}J^*$. The coefficients
are signed. Positive and negative terms can decay at different rates even
though the training residual norm is nonincreasing.

\begin{corollary}[A sufficient two-rate envelope]\label{cor:rates}
Choose $0\le s<f$. Let $A$ be the sum of positive $c_j$ with
$\lambda_j\le s$, let $B_-$ be the sum of $|c_j|$ over negative
$c_j$ with $\lambda_j\ge f$, and let $D_-$ be the sum of $|c_j|$
over all remaining negative coefficients. Then
\[
 \Psi(\tau)\ge A e^{-s\tau}-B_-e^{-f\tau}-D_-.
\]
Consequently if the right side exceeds $\rho$ on a nontrivial late
interval, and $\Psi<-\rho$ on a nontrivial earlier interval,
$\Phi_I(\theta_t)$ has a decreasing interval followed by an increasing
interval under Theorem~\ref{thm:comparison}.
\end{corollary}
\begin{proof}
Bound retained positive terms below by $e^{-s\tau}$ times their total
coefficient, the fast negative terms below by $-e^{-f\tau}B_-$,
and the remaining negative terms below by $-D_-$. Discard other positive
terms. Apply \eqref{eq:rho} and integrate the strict signs almost everywhere.
\end{proof}
The full spectral sum may give a much sharper certificate than this envelope;
no claim that two literal eigenmodes suffice is required.

\begin{theorem}[Crossing localization from an entry-state response]\label{thm:localization}
Suppose the comparison error is uniformly at most $\rho$ on $[0,H]$.
Assume the explicitly computed $\Psi$ has $0<\tau_*<H$, and for some
$\kappa>0$,
\[
 \Psi(\tau)\le-\kappa(\tau_*-\tau)\quad(\tau<\tau_*),\qquad
 \Psi(\tau)\ge\kappa(\tau-\tau_*)\quad(\tau>\tau_*).
\]
If $\rho/\kappa<\min(\tau_*,H-\tau_*)$, every minimum of
$\Phi_I(\theta_t)$ on $[a,a+H]$ is interior and is within
$\rho/\kappa$ of $a+\tau_*$. The function is strictly decreasing before
this bracket and strictly increasing after it. A continuous zero of $Q_G$
is not asserted.

If the exact full-probe discrepancy satisfies
$|\dot E_{\rm full}-Q_G|\le d$ on the same interval, replace $\rho$ by
$\rho+d$ to obtain the corresponding conclusion for $E_{\rm full}$.
\end{theorem}
\begin{proof}
For $\tau<\tau_*-\rho/\kappa$, the upper bound for $Q_G$ is strictly
negative almost everywhere. For $\tau>\tau_*+\rho/\kappa$, its lower
bound is strictly positive. Since $\Phi_I(\theta_t)$ is absolutely
continuous, integration proves strict monotonicity on the outer intervals.
Continuity and compactness give a minimum, and those intervals exclude any
minimizer outside the bracket. The proof for $E_{\rm full}$ uses the same
argument after adding $d$.
\end{proof}
In contrast with assuming the actual $Q_G$ crossing, the reference response
and its slopes are calculated from one state. Reachability of that state and
a sufficiently small error remain separate questions. The full discrepancy
$d$ must include inactive strong neurons, the active complement and their
velocities, and the difference between $Q_G$ and the exact reduced derivative.
A forward residual bound alone does not supply $d$.

\section{A completely explicit drift bound, including gate changes}
The preceding error can be bounded from a single state without postulating a
trajectory regularity constant. The following deliberately conservative bound
is useful for stating exactly what has and has not been certified.
Let $X_* = \max_n\|x_n\|$, $P=\|\theta_a\|$,
$R_* = \sqrt{H\mathcal L(a)}$, and $V=\|v\|$.
For each neuron define the set of possibly changing sample gates
\[
 \mathcal B_i=\{x_n:|u_i(a)^Tx_n|\le R_*\|x_n\|\},\qquad
 \Gamma^2=\sum_i\|w_i(a)\|^2\E_n[\|x\|^2\mathbf1_{\mathcal B_i}].
\]
Set
\begin{align*}
 D_M&=PR_*+\tfrac12R_*^2,&
 D_J&=\sqrt2 X_*R_*+\Gamma,\\
 D_g&=V\|\ell_a\|R_*+V^2D_M(P+R_*),& j&=\|J_a\|,\\
 D_C&=D_J\|g_a\|+(j+D_J)D_g,&
 D_K&=(2j+D_J)D_J.
\end{align*}

\begin{proposition}[Primitive finite-window error]\label{prop:primitive}
Throughout $[a,a+H]$, these constants bound respectively the block change,
Jacobian change, potential-gradient change, $C$ change and $K$ change.
In particular the comparison theorem holds with
\begin{equation}\label{eq:primitive}
 \rho_{\rm prim}=\sqrt{2\mathcal L(a)}
       \big[D_C+H D_K\|C_a\|_n\big].
\end{equation}
No probability claim or entry-state assertion is included.
\end{proposition}
\begin{proof}
Loss dissipation and Cauchy--Schwarz imply
\[
 \|\theta_t-\theta_a\|\le\int_a^t\|\dot\theta_s\|ds
 \le\sqrt{(t-a)(\mathcal L(a)-\mathcal L(t))}\le R_*.
\]
For a fixed block, expanding $W_I^TU_I$ gives
$\|\Delta M\|\le PR_*+R_*^2/2$.
A changed selector can only occur on $\mathcal B_i$. Split the Jacobian
change into its smooth weight/activation change and its changed-selector
part, attaching the latter to $w_i(a)$. The squared Hilbert--Schmidt norm
of the smooth part is at most
$X_*^2(2\|\Delta U\|_F^2+\|\Delta W\|_F^2)\le2X_*^2R_*^2$.
The selector part has Hilbert--Schmidt norm at most $\Gamma$.
The operator norm is no larger; this proves $D_J$.

Next $\|\ell_t-\ell_a\|=|Y_t-Y_a|\le V D_M$.
The explicit formulas for $g$ in Proposition~\ref{prop:potential} give
$\|g_t-g_a\|\le V\|\ell_a\|R_*+V^2D_M(P+R_*)$.
Expanding $J_tg_t-J_ag_a$ and $J_tJ_t^*-J_aJ_a^*$ gives $D_C,D_K$.
Apply Theorem~\ref{thm:comparison} and $\|e_a\|_n=\sqrt{2\mathcal L(a)}$.
\end{proof}
If every initial training margin exceeds $R_*\|x\|$, then $\Gamma=0$.
For bounded state norms, fixed $H$ and fixed positive margins, writing
$\epsilon_e=\|e_a\|_n$ gives $R_*=O(\epsilon_e)$,
$D_J,D_g,D_C,D_K=O(\epsilon_e)$, and therefore
$\rho_{\rm prim}=O(\epsilon_e^2)$. Thus an order-$\epsilon_e$ signed
spectral crossing is stable to this bound in such a small-residual regime.
This scaling observation does not prove that canonical random initialization
reaches either the spectral condition or the required margins.

\section{Checks on existing canonical states and the remaining inequality}
The reanalysis uses the saved canonical seed-$0$ trajectory with $h=200$,
$N=4000$, means $3,2$, noise standard deviation $0.15$, and the stored strong
mask selected at time $4$. The fixed probe is $-85^\circ$.
No training was rerun. The saved states come from Euler integration;
all numbers below are retrospective floating-point diagnostics, not validated
continuous-time bounds. The scripts reconstruct the original seeded dataset.

At $t=3.60$, equation~\eqref{eq:square} has the following terms:
\begin{center}
\begin{tabular}{lr}\toprule
Term & Value\\\midrule
Positive square & $+0.0036797850$\\
Complement correlation & $-0.0006266230$\\
Target correlation & $-0.0042214838$\\
Input--output mismatch contribution & $+0.0012836054$\\
Gate-weighted $x_1$ correction & $+0.0017669433$\\\midrule
Total $\mathcal S_2$ & $+0.0018822269$\\\bottomrule
\end{tabular}
\end{center}
The target term alone exceeds the square. The mismatch and $x_1$ corrections
are essential to this decomposition's sign. The normalized alignment of
$f_1-x_1$ with $H_I$ is approximately $0.485$, rather than a consequence
of an unweighted weak residual norm.

The spectral reference in \eqref{eq:spectrum}, computed separately at each
listed entry time, gives these interpolated negative-to-positive crossings:
\begin{center}
\begin{tabular}{rrrr}\toprule
Entry $a$ & $Q_G(a)$ & Reference crossing $a+\tau_*$ & $\Psi(0.5)$\\\midrule
2.9 & $-0.00724618$ & 3.813095 & $-0.00136129$\\
3.0 & $-0.00555350$ & 3.768320 & $-0.00077795$\\
3.2 & $-0.00281777$ & 3.704165 & $-0.00000939$\\
3.4 & $-0.00105117$ & 3.654520 & $+0.00043947$\\
3.5 & $-0.00045939$ & 3.630368 & $+0.00059510$\\\bottomrule
\end{tabular}
\end{center}
The supplied audit gives the actual-statistic crossing on saved states as
$3.60182403$, by interpolation. The frozen reference is therefore informative
but not exact. Starting at $a=3.4$, linear decomposition by the support of the
\emph{initial residual} gives at $\tau=0.5$ a strong-cluster contribution
$-0.0009715782$ and a weak-cluster contribution $+0.0014110470$.
Both contributions evolve under the full coupled reference kernel. This is
not an ablation removing the weak training cluster and is not the same
partition as the instantaneous clusterwise source in the supplied audit.

\paragraph{What the current bound actually certifies.}
For the diagnostic state $a=3.4$, $H=0.5$, the primitive formulas give
\[
 \mathcal L(a)=0.0041220000,\qquad R_*=0.04539824,
 \qquad \rho_{\rm prim}=0.40288812.
\]
The required positive-witness inequality would be
\begin{equation}\label{eq:failed}
 \rho < \Psi(0.5)=0.0004394687840.
\end{equation}
It is not satisfied by the primitive bound. This failure does not disprove
a useful finite-window theorem. In particular it does not show that the
actual signed discrepancy is large or adverse.

\paragraph{The sharper remaining target.}
For the late witness one only needs
\begin{equation}\label{eq:remaining}
\begin{split}
 &\ip{C_{a+H}-C_a}{e_{a+H}}\\
 &\quad-\int_0^H\ip{e^{-K_a(H-s)}C_a}{(K_{a+s}-K_a)e_{a+s}}\,ds
 >-\Psi(H)
\end{split}
\end{equation}
on a neighborhood of a late time (almost everywhere), together with an early
negative interval. A proved absolute bound satisfying \eqref{eq:failed} is
sufficient but not necessary. No estimate in the existing specialization or
loss-dissipation results establishes \eqref{eq:remaining} on a reached
canonical window. The source square does not fill this gap either.


\paragraph{A smaller signed budget on the saved states.}
At $a=3.4$, $H=0.5$, retrospective evaluation of
Corollary~\ref{cor:directional} gives
\begin{align*}
 (g_{a+H}-g_a)^Tb_H&=+0.0003228852,\\
 g_a^T\Delta J_H^*e_{a+H}&=+0.0001586021,\\
 P_H&=+0.0004814873,\qquad
 B_H^{\rm sampled}=0.0003116022.
\end{align*}
The last quantity is a trapezoidal sum of the pointwise norm products, not a
rigorous upper bound on the continuous-time integral. If one could prove
$P_H\ge0$ and an upper bound as small as this sampled $B_H$, the positive
margin left by \eqref{eq:directional-budget} would be about $0.0001278666$.
This shows why taking absolute values of the endpoint drift is counterproductive:
the corresponding sampled absolute budget is $0.0007930895$, which exceeds
$\Psi(H)$. The signed Duhamel reconstruction differs from the directly
evaluated discrepancy by $2.35\cdot10^{-7}$; this includes quadrature and
Euler-trajectory error and is not a certified error bound.

This sharper calculation identifies a quantitatively plausible target but
still proves no uniform inequality on the true initialized flow. Validating a
fixed-seed canonical trajectory would be a distinct, computer-assisted milestone;
it would not supply a high-probability iid theorem or license treating Euler
snapshots as exact gradient-flow states.

\paragraph{Scope of the saved-state comparison.}
The new results isolate a route from a reached ID state to a crossing without
assuming the crossing itself. They do not prove that iid isotropic initialization
reaches a favorable spectral state or that its signed drift is favorable.
For the comparison route, the coordinate-specific full-probe residual and velocity
must also satisfy the transfer budget on the same window. The separate local
certificate below instead verifies the full derivative and snapshots directly. No common noisy
high-probability parameter region or canonical timing guarantee is claimed.
The supplied noiseless countermodel remains compatible with every statement
here: its actual statistic need not satisfy a favorable spectral certificate.
The existing finite-sample theorem and cleaned appendix are unchanged.
\section{Static sign certificates with all training-gate changes}
\label{sec:static}
This section concerns parameter neighborhoods, not their reachability from
initialization. Its dataset consists of the exact binary numbers in the
accompanying \texttt{canonical\_witness\_states.npz}. Rounding a Gaussian
sample is not an iid Gaussian theorem.

Let $y$ be a parameter center, $R>0$, $P=\|y\|$, $\bar P=P+R$, and
$D=PR+R^2/2$. A cluster index set $C$ uses the full-dataset normalization
$\|z\|_C^2=N^{-1}\sum_{n\in C}\|z_n\|^2$. Put
$\lambda_C\ge\|N^{-1}\sum_{n\in C}x_nx_n^T\|$.
Write $F_C$ for that cluster's contribution to the full gradient field,
using the \emph{full-network residual}; $F=F_1+F_2$.
Let $F_C^0$ denote its polynomial extension with the training gates fixed
at $y$. No fixed-gate hypothesis is imposed on the actual field in the ball.

For either potential used here, set $g=\nabla\Phi$, $G_0=\|g(y)\|$.
For the empirical statistic use $\Phi=\Phi_I$ from
Proposition~\ref{prop:potential}. For the full derivative, fix the actual
probe-active set at $y$ and use
$\Phi(\theta)=\frac12\|M_{\rm act}(\theta)\xi-\xi\|^2$.
This second choice is valid throughout the ball if
$\min_i|u_i(y)^T\xi|>R\|\xi\|$.
In either case let $v$ be its input vector, $I$ its fixed index set,
$P_I=\|(u_i,w_i)_{i\in I}(y)\|$, and $\ell_y$ its output gradient.
Define
\begin{align*}
 L_C&=\lambda_C\bar P^2+
 \sqrt{\lambda_C}\big(\|e(y)\|_C+\sqrt{\lambda_C}D\big),\\
 H_g&=\|v\|^2(P_I+R)^2+
 \|v\|\big(\|\ell_y\|+\|v\|(P_IR+R^2/2)\big).
\end{align*}

Here are explicit bounds for the training-gate correction. Put
$q_{ni}=u_i(y)^Tx_n$, $s_n=\|x_n\|$, $e_n=e(y;x_n)$, and
\begin{align*}
 d_{ni}&=(Rs_n-|q_{ni}|)_+,&
 p_{ni}&=\mathbf1_{|q_{ni}|\le Rs_n},\\
 b_{ni}&=|w_i(y)^Te_n|+R\|e_n\|
                 +(\|w_i(y)\|+R)Ds_n,\\
 j_i^u&=N^{-1}\sum_{n\in C}p_{ni}b_{ni}s_n,&
 j_i^w&=N^{-1}\sum_{n\in C}d_{ni}(\|e_n\|+Ds_n),\\
 z_n&=\sum_i(\|w_i(y)\|+R)d_{ni},&
 B_C&=\Big[\sum_i((j_i^u)^2+(j_i^w)^2)\Big]^{1/2}
                  +\sqrt{\lambda_C}\bar P\|z\|_C.
\end{align*}
Every neuron is retained in these quantities.

\begin{proposition}[A static neighborhood enclosure]\label{prop:static}
For any compatible training selectors and $\|\theta-y\|\le R$,
\begin{equation}\label{eq:static}
 \left|g(\theta)^TF_C(\theta)-g(y)^TF_C(y)\right|
 \le R\left[H_g(\|F_C(y)\|+L_CR)+G_0L_C\right]
                   +(G_0+H_gR)B_C.
\end{equation}
This applies to the total derivative and separately to each cluster's
signed contribution. Cluster contributions always use the same full residual.
\end{proposition}
\begin{proof}
The Jacobian norm on cluster $C$ is at most
$\sqrt{\lambda_C}\|\theta\|$: apply Cauchy--Schwarz across neuron and
input/output parameter coordinates before taking the empirical average.
Integration along a parameter segment gives
$\|f_\theta-f_y\|_C\le\sqrt{\lambda_C}D$ and the pointwise bound
$\|f_\theta(x_n)-f_y(x_n)\|\le Ds_n$.
The same bounds hold for the fixed-mask polynomial extension.
Its field derivative is $-J_C^*J_C+H_{e,C}$, where the neuronwise
residual Hessian has norm at most
$\sqrt{\lambda_C}\|e\|_C$. Thus
$\|DF_C^0\|\le L_C$. Differentiating either potential gives
$\|Dg\|\le H_g$ on the ball. Adding and subtracting
$g(y)^TF_C^0(\theta)$ proves the first term of \eqref{eq:static}.

A changed training activation differs from its fixed-mask extension by at
most $d_{ni}$, and its selector can differ only when $p_{ni}=1$.
The pointwise residual and weight bounds give
$|w_i(\theta)^Te_\theta(x_n)|\le b_{ni}$. The direct input- and
output-gradient changes are therefore bounded by $j_i^u,j_i^w$.
The change of residual between the actual network and the fixed-mask
extension has norm at most $z_n$. Applying the fixed-mask adjoint to
this residual costs at most $\sqrt{\lambda_C}\bar P\|z\|_C$.
Consequently $\|F_C-F_C^0\|\le B_C$. Multiply by
$\|g(\theta)\|\le G_0+H_gR$ to finish. Endpoint partial selectors obey
the same bounds.
\end{proof}

\paragraph{Verified static values.}
The two centers are RK4 reference states at $3.4$ and $3.9$, generated with
step $1/5000$, the original canonical dataset and the exact archived initial
parameter values. The cohort is the original offline strong cohort of
$71$ neurons. The bounds below hold on balls of radius $3.5\cdot10^{-5}$.
Entries in the table are outward rounded and multiplied by $10^3$.
\begin{center}
\begin{tabular}{llrr}
\toprule
Statistic & contribution & center at $3.4$ & center at $3.9$\\
\midrule
$Q_G$ & total & $[-1.205,-0.894]$ & $[0.677,0.982]$\\
$Q_G$ & strong cluster & $[-2.479,-2.198]$ & $[-1.224,-0.944]$\\
$Q_G$ & weak cluster & $[1.211,1.367]$ & $[1.835,1.992]$\\
$\dot E$ & total & $[-1.439,-1.156]$ & $[0.449,0.728]$\\
$\dot E$ & strong cluster & $[-2.522,-2.266]$ & $[-1.234,-0.978]$\\
$\dot E$ & weak cluster & $[1.025,1.168]$ & $[1.622,1.766]$\\
\bottomrule
\end{tabular}
\end{center}
Thus weak-cluster forcing opposes the strong-cluster improvement at both
states and dominates it at the later state. This is a signed decomposition
of the actual field, not a retraining ablation or an assumed positive source.
The full-error computation includes all probe-active neurons directly;
it does not infer its sign from a small forward reduction residual.

On both balls the cohort lies in the positive strong cone: input angles
are below $15^\circ$, output angles below $20^\circ$. Its ungated
$m_{11}$ is greater than $0.9322$ and $0.9340$, respectively.
These are coarse learned-block geometry statements, not fine locking or
an initialization-to-specialization theorem. The full-field speed is at most
$0.385$ and $0.309$, respectively.

The accompanying script evaluates \eqref{eq:static} with outward
midpoint-radius bounds, dot-product rounding bounds, and interval
trigonometry. All center training masks are resolved by intervals; possible
changes anywhere in each ball are charged by $B_C$. The probe masks stay
fixed because their center margins exceed $3.7575\cdot10^{-5}$ and
$5.6601\cdot10^{-5}$. Exact-rational corner tests independently check
$553$ arithmetic enclosures. The arithmetic contract is IEEE
round-to-nearest with standard dot accumulation and no unsafe reassociation;
these tests are not a proof of initialized reachability.

\input{CANONICAL_ENCLOSURE_METHOD.tex}
\section{A local noisy many-neuron reversal certificate}
\label{sec:local-result}
The following result is a conditional entry theorem. Its hypothesis specifies
a neighborhood of an ID-trained state; it does not assume an OOD derivative
sign or a crossing. The local dynamics and the OOD conclusions are certified
below. Initialization-to-entry is a separate unresolved inequality.

Let $X$ be the supplied canonical binary dataset, $h=200$, and
$\widehat\theta_a$ the supplied numerical center indexed by $a=3.4$.
The loss uses all $N=4000$ samples. The fixed cohort $I$ is the original
$71$-neuron offline strong cohort. Let
\[
 \xi_\alpha=(\cos\alpha,\sin\alpha),\qquad
 \mathcal I=\{\alpha:|\alpha+17\pi/36|\le10^{-6}\}.
\]
The sector has angular width $2\cdot10^{-6}$ radians.

\begin{theorem}[Reversal from a certified learned-state neighborhood]
\label{thm:local-certificate}
Assume a compatible gradient-flow solution satisfies
\begin{equation}\label{eq:entry-missing}
 \|\theta(3.4)-\widehat\theta_{3.4}\|_2\le10^{-6}.
\end{equation}
Then, uniformly for $\alpha\in\mathcal I$,
\begin{align*}
 E(\xi_\alpha,3.4)-E(\xi_\alpha,3.7)&>1.2\cdot10^{-4},\\
 E(\xi_\alpha,3.9)-E(\xi_\alpha,3.7)&>4\cdot10^{-5}.
\end{align*}
In particular the full-network error attains an interior minimum on
$[3.4,3.9]$. At entry the fixed cohort has $m_{11}>0.9322$, input angles
below $15^\circ$, and output angles below $20^\circ$.
For the central probe, $Q_G<0$ at entry and $Q_G>0$ at $3.9$ for every
compatible training selector; the full-error rate has the same strict signs.
More quantitatively, for every probe in $\mathcal I$,
$\dot E<-0.00115$ almost everywhere on $[3.4,3.40001]$, and
$\dot E>0.000444$ almost everywhere on $[3.9,3.90001]$.
No continuous or unique zero of $Q_G$ is asserted.
\end{theorem}
\begin{proof}[Proof and computational obligations]
The local enclosure uses $2500$ cubic-reference steps of length $1/5000$.
At each step, outward bounds for the uniform reference defect, smooth
one-sided rate, and adverse gate jumps are inserted into
Corollary~\ref{cor:tube-step}. The input radius is $10^{-6}$, rounded upward for the scalar replay.
The certified radii at $3.7$ and $3.9$ are below
$1.199\cdot10^{-5}$ and $2.217\cdot10^{-5}$, respectively.
Every strict tube-invariance inequality must pass; the complete list is
\texttt{LOCAL\_FLOW\_ENCLOSURE.csv}, with its terminal report in
\texttt{LOCAL\_FLOW\_CERTIFICATE.json}. These computations certify that
the flow is within $3.5\cdot10^{-5}$ of the reference snapshots at $3.7$
and $3.9$. Section~\ref{sec:static} proves the static rate and geometry
bounds used here. The computational proof is conditional on the arithmetic
contract stated there; source code and exact binary inputs are supplied.

For completeness, the independent snapshot verification gives
\begin{align*}
 E(\xi_0,3.4)&\in[0.48314414,0.48315477],\\
 E(\xi_0,3.7)&\in[0.48298853,0.48299930],\\
 E(\xi_0,3.9)&\in[0.48306002,0.48307086],
 \qquad \xi_0=\xi_{-17\pi/36}.
\end{align*}
The displayed intervals are enlarged from the full-precision certified
values. They follow by Taylor's theorem for the full probe-error potential,
with error bound $\|g(y)\|R+H_gR^2/2$; all probe gates are constant
on the respective static balls. The first and second central gaps exceed
$0.00014484$ and $0.00006072$.

For any parameter vector in these balls,
$\sum_i\|w_i\|\|u_i\|\le\|\theta\|^2/2$.
Thus the full error on the unit circle is Lipschitz in the probe with
constant at most $(1+\|\theta\|^2/2)^2<9.084$.
Moving the probe by at most $10^{-6}$ in angle changes a difference of
two errors by less than $1.8168\cdot10^{-5}$. This proves the two
claimed uniform gaps. Continuity gives an interior minimum. The mass,
cone, and central derivative conclusions are the corresponding static
certificates. The terminal enclosure is smaller than $3\cdot10^{-5}$.
The speed bounds $0.385$ and $0.309$ ensure that extending either
endpoint forward by $10^{-5}$ stays in its radius-$3.5\cdot10^{-5}$
static ball, by a first-exit argument. Across the probe sector the masks
remain unchanged: the additional score movement is at most
$2.01\cdot10^{-6}$, below the unused central probe margin. On either ball
\[
 |\dot E(\xi)-\dot E(\xi_0)|
 \le2(1+P^2/2)P\|F\|\,\|\xi-\xi_0\|
 <5\cdot10^{-6},\qquad P<2.01.
\]
Subtract this cost from the static rate margins to obtain the asserted
uniform decreasing and increasing intervals. This also extends the
compatible flow through the short terminal interval.
\end{proof}

\paragraph{What the signed mechanism says.}
At both endpoints the strong-cluster contribution to the full OOD rate is
negative and the weak-cluster contribution is positive. The latter dominates
at $3.9$. Moreover its late lower bound exceeds its early upper bound, while
the strong-cluster contribution becomes less negative, by strict certified
endpoint margins. These are contributions to the \emph{same} gradient field
and full-network residual, not results of removing a cluster and retraining.
They explain the local reversal by signed competition between the two
training clusters. They do not prove that specialization alone is sufficient
or that it causes reversal for every specialized state.

\paragraph{The remaining initialized obligation.}
The certificate starts from the small neighborhood in
\eqref{eq:entry-missing}. It does not prove that the canonical initialized
flow reaches it. In particular, restarting a numerical or exact flow at
$\widehat\theta_{3.4}$ is not the same as certifying the path from the
archived seed. An initialized canonical theorem still requires
\eqref{eq:entry-missing}, or another proved entry enclosure for which the
local bounds close. The high-probability iid theorem requires an additional
probability argument and is not implied by this fixed-dataset result.

\end{document}
